\chapter*{Preface}

\vspace{0.5em}
Large language models originated from natural language processing, but they have undoubtedly become one of the most revolutionary technological advancements in the field of artificial intelligence in recent years. An important insight brought by large language models is that knowledge of the world and languages can be acquired through large-scale language modeling tasks, and in this way, we can create a universal model that handles diverse problems. This discovery has profoundly impacted the research methodologies in natural language processing and many related disciplines. We have shifted from training specialized systems from scratch using a large amount of labeled data to a new paradigm of using large-scale pre-training to obtain foundation models, which are then fine-tuned, aligned, and prompted.

This book aims to outline the basic concepts of large language models and introduce the related techniques. As the title suggests, the book focuses more on the foundational aspects of large language models rather than providing comprehensive coverage of all cutting-edge methods. The book consists of six chapters:

\begin{itemize}
\item Chapter 1 introduces the basics of pre-training. This is the foundation of large language models, and common pre-training methods and model architectures will be discussed here.
\item Chapter 2 introduces generative models, which are the large language models we commonly refer to today. After presenting the basic process of building these models, we will also explore how to scale up model training and handle long texts.
\item Chapter 3 introduces prompting methods for large language models. We will discuss various prompting strategies, along with more advanced methods such as chain-of-thought reasoning and automatic prompt design.
\item Chapter 4 introduces alignment methods for large language models. We will focus on instruction fine-tuning and alignment based on human feedback.
\item Chapter 5 introduces inference methods for large language models. We will discuss various decoding algorithms, acceleration methods, and the inference-time scaling issue.
\item Chapter 6 introduces methods for improving the reasoning ability of large language models. We will discuss how training can strengthen reasoning, and how additional computation during inference can help models solve complex problems.
\end{itemize}

If readers have some background in machine learning and natural language processing, along with a certain understanding of neural networks like Transformers, reading this book will be quite easy. However, even without this prior knowledge, it is still perfectly fine, as we have made the content of each chapter as self-contained as possible, ensuring that readers will not be burdened with too much reading difficulty.

The content presented here is part of a comprehensive introductory resource on neural networks and large language models in natural language processing. For readers who wish to learn more about background topics, such as sequence modeling and attention mechanisms, you can visit \url{https://github.com/NiuTrans/NLPBook} or \url{https://niutrans.github.io/NLPBook} for further information.

We would like to thank the students in our laboratory and all our friends who have shared with us their views on large language models and helped with corrections of errors in writing. In particular, we wish to thank Weiqiao Shan, Yongyu Mu, Chenglong Wang, Kaiyan Chang, Yuchun Fan, Hang Zhou, Chuanhao Lv, Xinyu Liu, Tao Zhou, Huiwen Bao, Tong Zheng, Junhao Ruan, Yingfeng Luo, Yuzhang Wu, Yifu Huo, Xiaoqian Liu, Yucheng Qiao, Junxiang Zhang, Miaohe Niu, Junxin Wang, Xin Liu, Xinyi Xu, Dan Chen, Xihan Yang, and Jingqi Xiao.

\vspace{3em}
\begin{flushright}
Tong Xiao and Jingbo Zhu

\textit{\today}
\end{flushright}
